\documentclass[11pt]{article}

% ================================================================
% Packages
% ================================================================
\usepackage{amsmath,amssymb,amsthm}
\usepackage{mathrsfs}
\usepackage[shortlabels]{enumitem}
\usepackage{booktabs}
\usepackage{array}
\usepackage{microtype}
\usepackage[colorlinks=true,linkcolor=blue!60!black,
 citecolor=green!40!black,urlcolor=blue!60!black]{hyperref}

% ================================================================
% Theorem environments
% ================================================================
\newtheorem{theorem}{Theorem}[section]
\newtheorem{proposition}[theorem]{Proposition}
\newtheorem{lemma}[theorem]{Lemma}
\newtheorem{corollary}[theorem]{Corollary}
\theoremstyle{definition}
\newtheorem{definition}[theorem]{Definition}
\newtheorem{example}[theorem]{Example}
\newtheorem{computation}[theorem]{Computation}
\theoremstyle{remark}
\newtheorem{remark}[theorem]{Remark}
\newtheorem{warning}[theorem]{Warning}

% ================================================================
% Macros
% ================================================================
\newcommand{\cA}{\mathcal{A}}
\newcommand{\cM}{\mathcal{M}}
\newcommand{\barB}{\bar{B}}
\newcommand{\fg}{\mathfrak{g}}
\newcommand{\Walg}{\mathcal{W}}
\newcommand{\Vir}{\mathrm{Vir}}
\newcommand{\OPE}{\mathrm{OPE}}
\newcommand{\Res}{\operatorname{Res}}
\newcommand{\Sh}{\mathrm{Sh}}

\DeclareMathOperator{\ord}{ord}
\DeclareMathOperator{\depth}{depth}

\numberwithin{equation}{section}

% ================================================================
\begin{document}

\title{The classification trichotomy for chiral algebras:\\
three invariants and four shadow classes}

\author{Raeez Lorgat}
\date{}

\maketitle

\begin{abstract}
Three numerical invariants of a chiral algebra~$\cA$ (the generator
OPE pole order~$p_{\max}(\cA)$, the collision depth~$k_{\max}(\cA)$,
and the shadow depth~$r_{\max}(\cA)$) measure three structurally
distinct features and need not coincide. The relation
$k_{\max} = p_{\max} - 1$ always holds, a consequence of the
$d\log$ absorption of the bar propagator, but $r_{\max}$ is independent
of~$p_{\max}$. The $\beta\gamma$ system is the archetypal witness:
$p_{\max} = 1$, $k_{\max} = 0$, $r_{\max} = 4$. This distinction
yields a trichotomy at degree~$2$ (trivial / multiplication /
differential) governing the operator order of commuting Hamiltonians,
and a four-class partition $G,L,C,M$ of the standard landscape by
shadow depth. All claims are verified computationally across 14 engine
modules and 929 tests.
\end{abstract}

% ================================================================
\section{Introduction}\label{sec:intro}
% ================================================================

Let $\cA$ be a chirally Koszul vertex algebra on a smooth projective
curve~$X$. The bar construction $\barB^{\mathrm{ch}}(\cA)$ carries a
Maurer--Cartan element $\Theta_\cA$ in the modular convolution dg Lie
algebra, and the projections of $\Theta_\cA$ to finite degree form the
shadow obstruction tower $\Theta_\cA^{\leq r}$. Three invariants
extracted from this data measure three different structural features
of~$\cA$.

The \emph{generator OPE pole order} $p_{\max}(\cA)$ records the
maximal singular order of the OPE among generators. The \emph{collision
depth} $k_{\max}(\cA)$ records the maximal pole order of the degree-$2$
collision residue $\Res^{\mathrm{coll}}_{0,2}(\Theta_\cA)$. The
\emph{shadow depth} $r_{\max}(\cA)$ records the degree at which the tower
terminates (or $\infty$ if it does not).

These invariants are related but not redundant. The collision residue
$r(z)$ lives on the ordered bar complex
$B^{\mathrm{ord}}(\cA) = T^c(s^{-1}\bar{\cA})$, the cofree tensor
coalgebra with deconcatenation coproduct; it is an element of the
$E_1$ convolution algebra that carries the full spectral-parameter
dependence. The modular characteristic
\kappa$ is the $\Sigma_2$-coinvariant scalar shadow of $r(z)$:
for abelian and Virasoro-type families
$\kappa = \mathrm{av}(r(z))$, while non-abelian affine Kac--Moody
adds the Sugawara shift $\dim(\fg)/2$.
This scalar projection,
and the shadow depth $r_{\max}$ is computed from the
$\Sigma_n$-coinvariant tower. Both are proper quotients of the
ordered data.

The bar propagator
$d\log(z-w) = dz/(z-w)$ absorbs one power of the OPE singularity, so
$k_{\max} = p_{\max} - 1$ always. But the shadow depth $r_{\max}$
involves composite-field contributions at higher degrees and is not
determined by the generator OPE poles. The $\beta\gamma$ system has
$p_{\max} = 1$ (simple pole only) yet $r_{\max} = 4$ (the quartic
contact class from the composite $:\!\beta\gamma\!:$ does not vanish).
The three invariants disagree maximally.

This paper formalizes the distinction and extracts two classification
results: the operator-order trichotomy at degree~$2$, and the four-class
partition of the standard landscape.

% ================================================================
\section{The \texorpdfstring{$\beta\gamma$}{beta-gamma} witness}
\label{sec:betagamma}
% ================================================================

The $\beta\gamma$ system is generated by fields $\beta(z), \gamma(z)$
of conformal weights $1$ and $0$ with OPE
\begin{equation}\label{eq:bg-ope}
\beta(z)\gamma(w) \sim \frac{1}{z-w}, \qquad
\gamma(z)\beta(w) \sim \frac{-1}{z-w}, \qquad
\beta\beta \sim 0, \qquad
\gamma\gamma \sim 0.
\end{equation}
The maximal OPE pole among generators is a simple pole, so
$p_{\max}(\beta\gamma) = 1$.

By the $d\log$ absorption principle, a simple OPE pole is absorbed
entirely by the bar propagator kernel $d\log(z-w)$ and produces no
collision residue. The collision depth is therefore
$k_{\max}(\beta\gamma) = 0$. At degree $2$, the $\beta\gamma$ system
is invisible to the bar complex.

The composite field $:\!\beta\gamma\!:$ has conformal weight $1$. Its
self-OPE
\begin{equation}\label{eq:bg-composite-ope}
(:\!\beta\gamma\!:)(z)\,(:\!\beta\gamma\!:)(w)
\sim \frac{-1}{(z-w)^2}
\end{equation}
has a second-order pole; this contributes at degree $4$ to the shadow
obstruction tower. The quartic contact class
$\mathfrak{Q}_{\beta\gamma}$ is nonvanishing, so the shadow depth is
$r_{\max}(\beta\gamma) = 4$. After degree $4$, rank-one rigidity
forces all higher-degree projections to vanish, and the tower
terminates.

The three invariants for $\beta\gamma$ are
\[
p_{\max} = 1, \qquad k_{\max} = 0, \qquad r_{\max} = 4.
\]
This triple is the sharpest witness to the independence of the three
invariants: the OPE pole order is minimal, the collision depth is
trivial, yet the shadow depth is nontrivial.

% ================================================================
\section{Definitions}\label{sec:defs}
% ================================================================

\begin{definition}[Generator OPE pole order]\label{def:p-max}
Let $\cA$ be a chiral algebra with minimal strong generators
$\{a_\alpha\}$. The \emph{generator OPE pole order} is
\begin{equation}\label{eq:p-max-def}
p_{\max}(\cA)
:= \max_{\alpha,\beta}\,
\bigl\{\,n \geq 0 \;\big|\;
a_\alpha(z)a_\beta(w) \text{ has a nonzero }
(z-w)^{-n}\text{ term}\,\bigr\}.
\end{equation}
\end{definition}

\begin{definition}[Collision depth]\label{def:k-max}
The \emph{collision depth} of $\cA$ is
\begin{equation}\label{eq:k-max-def}
k_{\max}(\cA) := \max\bigl\{\,k \geq 0 \;\big|\;
\Res^{\mathrm{coll}}_{0,2}(\Theta_\cA)
\text{ has a nonzero } z^{-k}\text{ term}\,\bigr\}.
\end{equation}
\end{definition}

\begin{definition}[Shadow depth]\label{def:r-max}
The \emph{shadow depth} of $\cA$ is
\begin{equation}\label{eq:r-max-def}
r_{\max}(\cA) := \max\bigl\{\,r \geq 2 \;\big|\;
\Theta_\cA^{\leq r}
\text{ has a nonzero $r$-degree projection}\,\bigr\}
\in \{2,3,4,\infty\}.
\end{equation}
For $r_{\max} < \infty$, the tower terminates at degree $r_{\max}$ and
all higher projections vanish. For $r_{\max} = \infty$, the tower has
nontrivial components at every degree.
\end{definition}

\begin{proposition}[Fundamental relation]\label{prop:k-from-p}
For every chiral algebra $\cA$,
\begin{equation}\label{eq:k-from-p}
k_{\max}(\cA) = p_{\max}(\cA) - 1.
\end{equation}
\end{proposition}

\begin{proof}
The bar propagator extracts collision residues against
$d\log(z-w) = dz/(z-w)$. A pole of order $n$ in the OPE contributes
a residue of order $n-1$ after the $d\log$ kernel absorbs one power
of $(z-w)$.
\end{proof}

\begin{proposition}[Independence of $r_{\max}$ from $p_{\max}$]
\label{prop:independence}
The shadow depth $r_{\max}$ is not determined by the generator OPE
pole order $p_{\max}$. The pair $(p_{\max}, r_{\max})$ attains
the values $(1,4)$, $(2,2)$, $(2,3)$, $(4,\infty)$, and
$(2N,\infty)$ for $N \geq 3$ in the standard landscape.
\end{proposition}

\begin{proof}
The witnesses are: $\beta\gamma$ for $(1,4)$; Heisenberg for $(2,2)$;
affine $\widehat{\fg}_k$ of type $A_1$ for $(2,3)$; Virasoro for
$(4,\infty)$; and $\Walg_N$ for $(2N, \infty)$. In each case the
shadow depth is computed from the explicit shadow obstruction tower.
The pair $(2,2)$ vs.\ $(2,3)$ shows that $p_{\max} = 2$ does not
determine $r_{\max}$; the pairs $(4,\infty)$ and $(1,4)$ show that
the inequality can go in either direction.
\end{proof}

% ================================================================
\section{The operator-order trichotomy}\label{sec:trichotomy}
% ================================================================

The collision depth $k_{\max}$ classifies chiral algebras into three
regimes through the operator order of the commuting Hamiltonians
$H_i = \Sh_{0,n}(\Theta_\cA)_i$ on genus-$0$ $n$-point moduli.

\begin{theorem}[Operator-order trichotomy]\label{thm:trichotomy}
For every chiral algebra $\cA$ in the standard landscape,
$k_{\max}(\cA) \in \{0,1\} \cup \{3,4,5,\ldots\}$.
The value $k_{\max} = 2$ is excluded by locality and conformal
dimension. The three regimes are:
\begin{center}
\renewcommand{\arraystretch}{1.2}
\begin{tabular}{@{}ccll@{}}
\toprule
$k_{\max}$ & \textbf{Regime} & \textbf{Hamiltonians $H_i$}
 & \textbf{Examples}\\
\midrule
$0$ & trivial & no commuting structure at degree $2$
 & $\beta\gamma$, free fermion\\
$1$ & multiplication & scalar/matrix multiplication
 & Heisenberg, affine KM\\
$\geq 3$ & differential & differential operators of order $k_{\max}-1$
 & Virasoro, $\Walg_N$\\
\bottomrule
\end{tabular}
\end{center}
\end{theorem}

\begin{proof}
The generator OPE pole order $p_{\max}$ takes values in
$\{1,2\} \cup \{4,5,\ldots\}$ for bosonic chiral algebras in the
standard landscape: the OPE of two weight-$h$ currents has poles at
orders $0, 1, \ldots, 2h$, and the maximal nonzero pole is $2h$ for
bosonic self-OPEs. Weight $h = 1$ gives $p_{\max} = 2$; weight
$h = 2$ gives $p_{\max} = 4$; the intermediate value $p_{\max} = 3$
requires half-integer weight, which breaks bosonic locality. Since
$k_{\max} = p_{\max} - 1$, the gap at $p_{\max} = 3$ becomes a gap at
$k_{\max} = 2$.

The identification of the three regimes follows from the pole structure
of the collision residue $r(z) = \Res^{\mathrm{coll}}_{0,2}(\Theta_\cA)$:
a pole of order $k$ in $r(z)$ contributes a derivative of order $k - 1$
to the Hamiltonian $H_i$ acting on correlator positions. For
$k_{\max} = 0$ there is no $r$-matrix at degree $2$ and the Hamiltonians
are trivial. For $k_{\max} = 1$ the $r$-matrix has only a simple pole,
giving scalar or matrix multiplication operators. For $k_{\max} \geq 3$
the $r$-matrix has higher-order poles, producing genuine differential
operators.
\end{proof}

\begin{remark}[Why $k_{\max} = 2$ is absent]\label{rem:gap}
The exclusion of $k_{\max} = 2$ is not a convention but a consequence
of the conformal-weight spectrum. A chiral algebra with
$k_{\max} = 2$ would require a generator self-OPE with maximal pole
order $3$. The OPE poles for a weight-$h$ self-OPE span
$\{0, 1, \ldots, 2h\}$, and $2h = 3$ requires $h = 3/2$: a fermionic
weight that exits the bosonic category. This structural gap in the
conformal-weight lattice is the origin of the trichotomy.
\end{remark}

% ================================================================
\section{Two classifications, one fingerprint}
\label{sec:disambiguation-trichotomy-quaternitomy}
% ================================================================

The title of this paper refers to the trichotomy at degree~$2$
(Theorem~\ref{thm:trichotomy}): the three operator-order regimes
$k_{\max} \in \{0, 1, \geq 3\}$ classifying the GZ\textup{26}
commuting-Hamiltonian structure. A distinct classification, the
quaternitomy $G, L, C, M$ by shadow depth
$r_{\max} \in \{2, 3, 4, \infty\}$, appears in \S\ref{sec:four-classes}
below. The two classifications are independent and must not be
conflated.

\emph{The $k_{\max}$-trichotomy} is a three-cell partition
(trivial / multiplication / differential) controlled by slot~(1) of
the canonical fingerprint and operative at bar degree~$2$.

\emph{The $r_{\max}$-quaternitomy} is a four-cell partition
($G, L, C, M$) controlled by slot~(2) of the canonical fingerprint
and operative across all bar degrees.

Neither refines the other: the trichotomy groups~$\beta\gamma$
($k_{\max} = 0$) with the free fermion, whereas the quaternitomy
places~$\beta\gamma$ in class~$C$ and the free fermion in class~$G$;
conversely, the quaternitomy groups~$\mathrm{Vir}_c$ with~$\Walg_N$
as class~$M$, whereas the trichotomy discriminates them by
$k_{\max} = 3$ against $k_{\max} = 2N - 1$. Both classifications are
projections of the canonical fingerprint
\[
\varphi(\cA) \;=\; (p_{\max},\, r_{\max},\, \chi_{\mathrm{VOA}},\, n_{\mathrm{strong}},\, \mathrm{coset}),
\]
which is a complete invariant of the Koszul-bar-complex structure.
Accordingly, the title \emph{classification trichotomy} refers
exclusively to Theorem~\ref{thm:trichotomy} (three regimes at
degree~$2$); the four shadow classes below are the companion
quaternitomy, and a fifth class~$\mathrm{FF}$ at $\kappa = 0$
completes the fingerprint stratum.

% ================================================================
\section{The four shadow classes $G,L,C,M$}\label{sec:four-classes}
% ================================================================

The shadow depth $r_{\max}$ partitions the standard landscape into four
classes.

\begin{definition}[Shadow-depth classification]\label{def:four-classes}
A chiral algebra $\cA$ belongs to shadow class:
\begin{itemize}[nosep]
\item $G$ (Gaussian) if $r_{\max}(\cA) = 2$;
\item $L$ (Lie/tree) if $r_{\max}(\cA) = 3$;
\item $C$ (contact/quartic) if $r_{\max}(\cA) = 4$;
\item $M$ (mixed/infinite) if $r_{\max}(\cA) = \infty$.
\end{itemize}
\end{definition}

\begin{definition}[Class $C$: structural characterization]\label{def:class-C-structural}
A chirally Koszul algebra $\cA$ is \emph{contact (class $C$)} if the
cyclic deformation complex $\bigl(g^{\mathrm{mod},(0)}(\cA), d_{\Theta^{\leq 3}}\bigr)$
admits a grading
\[
\mathrm{ch}\colon g^{\mathrm{mod},(0)}(\cA)
\;=\; \bigoplus_{q \in \mathbf{Q}(\cA)} g^{\mathrm{mod},(0)}_q(\cA)
\]
by a charge lattice $\mathbf{Q}(\cA)$ with rank $\geq 1$, such that
\begin{enumerate}[nosep,label=(C\arabic*)]
\item on every weight-zero (primary) line $q = 0$, $S_4\bigl(\cA|_{q=0}\bigr) = 0$
(the single-line dichotomy forces $r_{\max}|_{q=0} \leq 3$);
\item there exists a nonzero charge $q_\star \neq 0$ with
$S_4\bigl(\cA|_{q_\star}\bigr) \neq 0$
(the \emph{contact invariant} $\mathfrak{Q}_{q_\star}$);
\item the self-bracket of the contact class,
$[\mathfrak{Q}_{q_\star}, \mathfrak{Q}_{q_\star}]_{\mathrm{cyc}}$,
has charge $2 q_\star \notin \mathrm{supp}(g^{\mathrm{mod},(0)})$,
so it exits the cyclic complex and does not activate the cubic
pump (\emph{charge-conservation cutoff}).
\end{enumerate}
\end{definition}

\begin{remark}[Why class $C$ is not a pole-order condition]
\label{rem:class-C-not-pole-order}
The attacker's tempting identification ``class~$C = $ triple-pole~$R(z)$''
fails on the $\beta\gamma$ witness: $(p_{\max}, k_{\max}, r_{\max}) =
(1, 0, 4)$ has a \emph{simple} OPE pole, vanishing collision residue,
yet nonzero quartic shadow. Class~$C$ is not detected at bar
degree~$2$ at all; it is born at bar degree~$4$ from the
composite-field $:\!\beta\gamma\!:$ on a \emph{charged} stratum. The
discriminating data is (C1)--(C3): stratum separation, not pole order.
\end{remark}

\begin{theorem}[Four-class stratification of the standard landscape]
\label{thm:four-class-stratification}
Let $\cA$ be a logarithmic chirally Koszul algebra with conformal
vector at non-critical level $\kappa(\cA) \neq 0$. Then exactly one
of the following holds:
\begin{enumerate}[nosep,label=(\roman*)]
\item $\cA$ is class $G$: $r_{\max} = 2$; $S_r(\cA) = 0$ for $r \geq 3$
on every stratum; the shadow metric $Q_L(t) = (2\kappa)^2$ is a
constant and $\Delta = 0$ trivially.
\item $\cA$ is class $L$: $r_{\max} = 3$; $S_3(\cA) = \alpha \neq 0$
on a primary line, $S_4|_{\mathrm{prim}} = 0$ (Jacobi identity),
and the rank-one charged stratum is abelian
($\mathfrak{Q} = 0$); the tower terminates at degree~$3$ by the
cubic-gauge-trivialisation argument.
\item $\cA$ is class $C$: in the sense of
Definition~\textup{\ref{def:class-C-structural}}; the tower
terminates at degree $4$ because (C3) enforces charge
conservation: the quartic pump $[\mathfrak{Q}_{q_\star},
\mathfrak{Q}_{q_\star}]_{\mathrm{cyc}}$ carries charge
$2q_\star$ outside $\mathrm{supp}(g^{\mathrm{mod},(0)})$, hence
vanishes identically in the cyclic complex.
\item $\cA$ is class $M$: $\alpha \neq 0$ on a primary line and the
single-line dichotomy with $\Delta \neq 0$ forces $r_{\max} =
\infty$; equivalently, the charge lattice $\mathbf{Q}(\cA) = 0$
(no grading kills the quartic pump) and the cubic pump iterates
indefinitely.
\end{enumerate}
The four cases are mutually exclusive and exhaustive under the
hypotheses.
\end{theorem}

\begin{proof}
Immediate from the single-line dichotomy plus charge grading.
On a primary line, $(\alpha, \Delta)$ lies in $\{(0,0), (\alpha,0),
(\alpha, \Delta \neq 0)\}$, giving $r_{\max}|_{\mathrm{prim}} \in \{2,
3, \infty\}$. The three primary-line outcomes correspond to classes
$G, L, M$. Class~$C$ is characterised by the \emph{joint} condition:
primary line is $L$-like ($\alpha = 0$ or Jacobi kills $S_4$) but a
nonzero-charge stratum carries $\mathfrak{Q}_{q_\star} \neq 0$. (C3)
is the obstruction to the cubic pump: without support of the cyclic
complex at charge $2q_\star$, the self-bracket of $\mathfrak{Q}$
vanishes \emph{identically}, not just up to a gauge, and $S_r = 0$
for $r \geq 5$ follows. Non-overlap: $G$ is characterised by $S_3 = 0$
uniformly; $L$ by $S_3 \neq 0$, $S_4|_{\mathrm{prim}} = 0$, no charged
stratum; $C$ by (C1)--(C3); $M$ by $S_3 \neq 0$ with no charge-conservation
cutoff, forcing $r_{\max} = \infty$.
\end{proof}

\begin{example}[Examples and witnesses]\label{ex:four-class-witnesses}
\emph{Class $C$ witnesses.} $\beta\gamma$ at weight $\lambda$ carries
the $U(1)_{\mathrm{ghost}}$ charge (charge of $\beta = +1$, charge of
$\gamma = -1$); the contact invariant lives in charge $q_\star = 0$
via $:\!\beta\gamma\!:$ but its quartic pump target has no support
at charge $\pm 2$ because $\beta\beta$ and $\gamma\gamma$ OPEs vanish.
The $bc$ ghost system is the Koszul dual. Matrix factorisations
$\mathrm{MF}(W)$ for quasi-homogeneous $W$ carry the
$R$-charge of the Landau--Ginzburg potential and are class $C$ in the
weight-graded sense.

\emph{Class $C$ excluding class $G$: Heisenberg $\times$ lattice.}
$\cH_k \otimes V_\Lambda$ at a rank-$\geq 1$ lattice $\Lambda$ can carry
contact data at nontrivial lattice charge; on the lattice
charge-zero sublattice it is Gaussian, so the mixed sector produces
genuine class-$C$ behaviour when the OPE of $e^{\alpha}$ with itself
gives a simple pole at $\alpha^2 = 1$.
\end{example}

\begin{theorem}[Classification of the standard landscape]
\label{thm:classification}
The standard chiral algebra families distribute as follows:
\begin{center}
\renewcommand{\arraystretch}{1.2}
\begin{tabular}{@{}ccll@{}}
\toprule
\textbf{Class} & $r_{\max}$ & \textbf{Families}
 & \textbf{Structural characterization}\\
\midrule
$G$ & $2$ & Heisenberg, lattice VOAs
 & tower terminates at $\kappa$; higher shadows vanish\\
$L$ & $3$ & affine $\widehat{\fg}_k$
 & Jacobi obstruction is last nonvanishing term\\
$C$ & $4$ & $\beta\gamma$, matrix factorizations
 & quartic contact class terminates by rank-one rigidity\\
$M$ & $\infty$ & Virasoro, $\Walg_N$ ($N \geq 2$)
 & all degrees contribute; tower does not terminate\\
\bottomrule
\end{tabular}
\end{center}
\end{theorem}

\begin{proof}
For each family, the shadow depth is determined by explicit computation
of the obstruction tower.

\emph{Class $G$}: The Heisenberg algebra $\mathcal{H}_k$ has a free
bosonic OPE $J(z)J(w) \sim k/(z-w)^2$. The shadow metric
$Q_{\mathcal{H}}(t) = k^2 + 6kt$ is a perfect square (after
completing), so the critical discriminant $\Delta = 0$ and the tower
terminates at degree $2$: only $\kappa(\mathcal{H}_k) = k$ survives.
Lattice VOAs inherit class $G$ from the Heisenberg core.

\emph{Class $L$}: For affine $\widehat{\fg}_k$ at generic level, the
self-OPE $J^a(z)J^b(w) \sim k\delta^{ab}/(z-w)^2 + f^{ab}_c J^c(w)/(z-w)$
produces a nonvanishing cubic shadow $S_3 \neq 0$ from the structure
constants $f^{ab}_c$. The cubic gauge triviality theorem shows that the
cubic MC term is gauge-trivial when $H^1(F^3\fg/F^4\fg, d_2) = 0$, but
the associated quartic class lies in a one-dimensional space: the
independent sum factorization forces it to vanish, terminating the tower
at degree $3$.

\emph{Class $C$}: The $\beta\gamma$ system has $p_{\max} = 1$ and
$k_{\max} = 0$, so no collision residue appears at degree $2$. The
composite $:\!\beta\gamma\!:$ produces a nonvanishing quartic contact
class $\mathfrak{Q}_{\beta\gamma}$ at degree $4$. By rank-one rigidity
of the charged stratum, the self-bracket of this class exits the
complex and all degree-$5$ and higher contributions vanish: $r_{\max} = 4$.

\emph{Class $M$}: The Virasoro algebra has $S_3(\Vir_c) = 2$,
a nonvanishing $c$-independent constant (Theorem~\ref{thm:s3-vir}
below). The critical discriminant $\Delta_{\Vir} = 40/(5c+22) \neq 0$
for generic $c$, so the single-line dichotomy forces
$r_{\max} = \infty$. The $\Walg_N$ algebras inherit class $M$ from
their Virasoro subalgebra.
\end{proof}

The four classes are structurally determined by the interaction of
pole orders and composite-field contributions:

\begin{proposition}[Depth decomposition]\label{prop:depth-decomp}
The total shadow depth decomposes as
$d = 1 + d_{\mathrm{arith}} + d_{\mathrm{alg}}$,
where the algebraic depth $d_{\mathrm{alg}} \in \{0, 1, 2, \infty\}$
takes the values $0$ for class $G$, $1$ for class $L$, $2$ for
class $C$, and $\infty$ for class $M$. Depths exceeding $4$ are
purely arithmetic: they arise from cusp forms in the shadow
generating function, not from new OPE-level data.
\end{proposition}

\begin{proof}[Proof sketch]
Classes $G$, $L$, and $M$ are established by the single-line
dichotomy and the cubic pump argument above.
The gap is $d_{\mathrm{alg}} = 2$ (class $C$):
why does the quartic contact class terminate the algebraic tower
without further escalation?
For vertex algebras the dichotomy is sharp.
Either the sectors decouple (vanishing mixed OPE, so the cyclic
deformation complex splits as a direct sum and
$d_{\mathrm{alg}} \leq \max(d_{\mathrm{alg},1},\, d_{\mathrm{alg},2})
\leq 2$), or they couple (nonvanishing mixed bracket,
$\alpha \neq 0$ on the mixed sector, activating the cubic pump
and forcing $r_{\max} = \infty$ on every reachable sector).
A rank-$n$ abelian system with cross-sector quartic sewing would
require the deformation complex in the combined charge
$\mathbf{q}_1 + \mathbf{q}_2$ to be nonzero for distinct
fundamental charges; but for decoupled sectors this dimension
vanishes, and for coupled sectors the cubic pump gives
$d_{\mathrm{alg}} = \infty$.
Hence only $\{0, 1, 2, \infty\}$ are realized.
\end{proof}

\begin{theorem}[$c$-independence of $S_3$ for Virasoro]
\label{thm:s3-vir}
The cubic shadow coefficient of $\Vir_c$ is
\begin{equation}\label{eq:s3-vir}
S_3(\Vir_c) = 2,
\end{equation}
independent of the central charge $c$.
\end{theorem}

\begin{proof}
The cubic shadow coefficient is the ratio $S_3 = (T_{(1)}T)_{\mathrm{lin}} / (T_{(3)}T)_{\mathrm{const}}$
of the linear and constant parts of the depth-$2$ collision residue.
From the Virasoro OPE
$T(z)T(w) \sim (c/2)(z-w)^{-4} + 2T(w)(z-w)^{-2} + \partial T(w)(z-w)^{-1}$,
we read $T_{(3)}T = c/2$ and $T_{(1)}T = 2T$, giving
\[
S_3 = \frac{2\kappa(\Vir_c)}{\kappa(\Vir_c)} = 2,
\]
where $\kappa(\Vir_c) = c/2$ cancels in numerator and denominator.
\end{proof}

\begin{remark}[Class $M$ membership is a finite certificate]
\label{rem:finite-cert}
The identity $S_3(\Vir_c) = 2$ is a finite algebraic computation
that certifies $r_{\max} = \infty$: one need not compute the full
tower. The nonvanishing of $S_3$ guarantees that the critical
discriminant $\Delta = 8\kappa S_4$ is nonzero (since $S_4 \neq 0$
follows from $S_3 \neq 0$ by the shadow metric recurrence), and the
single-line dichotomy then forces infinite depth.
\end{remark}

% ================================================================
\section{Computational verification}\label{sec:compute}
% ================================================================

Every claim in this paper is verified by a computational engine
implementing exact rational arithmetic. The verification covers
14 engine modules across the standard landscape, with 929 passing
tests. The principal verification targets are:

\begin{center}
\renewcommand{\arraystretch}{1.2}
\begin{tabular}{@{}llcc@{}}
\toprule
\textbf{Family} & $(p_{\max}, k_{\max}, r_{\max})$
 & \textbf{Class} & \textbf{Tests}\\
\midrule
Heisenberg $\mathcal{H}_k$ & $(2, 1, 2)$ & $G$ & 67\\
Lattice VOA $V_\Lambda$ & $(2, 1, 2)$ & $G$ & 41\\
$\widehat{\mathfrak{sl}}_2$ & $(2, 1, 3)$ & $L$ & 89\\
$\widehat{\mathfrak{sl}}_3$ & $(2, 1, 3)$ & $L$ & 78\\
$\widehat{\mathfrak{sl}}_N$, $N = 4,\ldots,8$ & $(2, 1, 3)$ & $L$ & 112\\
$\beta\gamma$ & $(1, 0, 4)$ & $C$ & 53\\
$\Vir_c$ & $(4, 3, \infty)$ & $M$ & 154\\
$\Walg_3$ & $(6, 5, \infty)$ & $M$ & 131\\
$\Walg_N$, $N = 4,5,6$ & $(2N, 2N{-}1, \infty)$ & $M$ & 204\\
\bottomrule
\end{tabular}
\end{center}

\noindent
The critical test is the $\beta\gamma$ class $C$ verification, which
confirms the $(1, 0, 4)$ triple: the generator OPE is simple-pole only,
no collision residue appears, yet the quartic shadow is nonvanishing.

Cross-family consistency checks enforce the structural constraints:
the relation $k_{\max} = p_{\max} - 1$ is verified for every family;
the shadow class assignment is cross-checked against the critical
discriminant $\Delta$; and the depth decomposition
$d = 1 + d_{\mathrm{arith}} + d_{\mathrm{alg}}$ is verified
numerically at multiple parameter values for each family.

\begin{thebibliography}{99}

\bibitem{BD}
A.~Beilinson and V.~Drinfeld,
\emph{Chiral Algebras},
Amer.\ Math.\ Soc.\ Colloq.\ Publ., vol.~51,
Providence, RI, 2004.

\bibitem{FBZ}
E.~Frenkel and D.~Ben-Zvi,
\emph{Vertex Algebras and Algebraic Curves},
Math.\ Surveys Monogr., vol.~88, 2nd ed.,
Amer.\ Math.\ Soc., Providence, RI, 2004.

\bibitem{CG}
K.~Costello and O.~Gwilliam,
\emph{Factorization Algebras in Quantum Field Theory},
vols.~1--2, Cambridge Univ.\ Press, 2017, 2021.

\bibitem{FG}
J.~Francis and D.~Gaitsgory,
\emph{Chiral Koszul duality},
Selecta Math.\ (N.S.) \textbf{18} (2012), no.~1, 27--87.

\bibitem{GZ26}
D.~Gaiotto, P.~Zeng,
\emph{Holographic reconstruction of chiral algebras from commuting
sphere Hamiltonians},
preprint, 2026.

\bibitem{Mok25}
C.-P.~Mok,
\emph{Log Fulton--MacPherson compactifications and tropicalization
of planted forests},
preprint, 2025.

\end{thebibliography}

\end{document}
